\chapter{Stability and Robustness Analysis}\label{ch:analysis}

The INDI loops cancel the plant, so their stability does not follow from the usual rate-loop tuning rules.
It rests on delays, filters and the actuator model.
A wrong model or a slow rotor can push a loop towards instability.
This chapter asks how far each loop is from that edge.
It measures the margins in flight simulation on the final configuration.
It then compares the measurements with linear models.
\Cref{ch:inner_loop,ch:outer_loop} describe the loops.

\section{How the loop gains are measured}\label{sec:an_method}

\subsection{Loop-break points}\label{sec:an_break}

A loop gain $L(s)$ is measured by opening the loop at one point and injecting a known signal.
The controller output $c$ and the excitation $d$ are added:
\begin{equation}\label{eq:an_inj}
  u = c + d .
\end{equation}
The sum $u$ goes to the plant and to the stored previous setpoint of the incremental law.
Every feedback path therefore closes through $u$: the gyro, the measured rotor speed and the modelled actuator.
The loop gain is the ratio of the signal that returns to the signal that was injected,
\begin{equation}\label{eq:an_L}
  L(j\omega) = -\frac{Y(j\omega)}{X(j\omega)} ,
\end{equation}
estimated from the logged signals.
Each loop has its own break point (\cref{tab:an_break}).

\begin{table}[tbp]
\caption{Loops, break points and injected signals.}\label{tab:an_break}
\centering
\small
\begin{tabularx}{\textwidth}{@{}L{3.6cm} L{4.6cm} Y@{}}
\toprule
Loop & Break point & Band of the multisine \\
\midrule
Rate (roll, pitch, yaw) & Torque setpoint, $u=c+d$ & 1 to 60~Hz \\
Attitude & Rate setpoint, output of the attitude loop & 0.2 to 20~Hz \\
Acceleration (INDI stage) & Output of the INDI stage (\texttt{pt1}) & 0.2 to 8~Hz \\
Velocity and position & Input of the INDI stage (\texttt{pt0}) & 0.2 to 8~Hz \\
\bottomrule
\end{tabularx}
\end{table}

The acceleration break closes every path through the plant input.
With a matched actuator model, the incremental terms cancel there.
The measured $L$ is then dominated by the outer loops.
The part that is specific to INDI shows above about 3~Hz.
The rate loop is broken at the torque, so its measurement contains the whole INDI law.

\subsection{Excitation and estimation}\label{sec:an_excite}

The excitation is a Schroeder multisine in hover.
Its frequency grid is the multiple of $1/T_p$, with 10 bins per decade.
It uses at least three periods, and the first one is discarded.
The amplitude is chosen in a pre-test: two amplitudes must agree within 1~dB and 5~deg.

The frequency response is estimated by a period-exact Fourier transform.
When the time base is not uniform, a least-squares fit on the true sample times is used instead.
Bins are kept where the coherence is at least 0.8, and at least 0.9 within a decade of the crossover.
Each run must also pass checks:
\begin{itemize}
  \item no fault and no landing during the run;
  \item the logged $u$ equals $c+d$;
  \item the logged rate equals the controller rate;
  \item the log has no holes beyond the allowed gap length.
\end{itemize}
Outer-loop cells were flown one at a time, because parallel flights made holes in the simulation log.

Every outer-loop cell is scored by two independent estimators, called key~1 and key~2.
They agree if the phase margins differ by at most 3~deg and the gain margins by at most 1~dB.
They must also agree on whether a gain margin exists.
A cell that fails this test is marked \emph{blocked}.

\section{Margin criteria}\label{sec:an_criteria}

The phase margin is $\mathrm{PM}=180^\circ+\angle L(j2\pi\fc)$ at the gain crossover $\fc$.
The gain margin is $\mathrm{GM}=-20\log_{10}|L|$ at the frequency where $\angle L=-180^\circ$.
The delay margin is the extra delay that brings the loop to the edge:
\begin{equation}\label{eq:an_dm}
  \mathrm{DM} = \frac{\mathrm{PM}}{360\,\fc} .
\end{equation}
The loops are accepted if all of the following hold.
\begin{itemize}
  \item $\mathrm{GM}\ge 6$~dB.
  \item $\mathrm{PM}\ge 45^\circ$.
  \item $\mathrm{DM}\ge 2$ control periods plus the telemetry latency.
\end{itemize}
The limits are applied to the worst of two runs, at controller rates of 800 and 400~Hz.
The lowest rate (about 267~Hz) is only reported.

The values are common margins for a loop that must stay stable when the model is wrong.
A gain margin of 6~dB tolerates a doubling of the effectiveness.
A phase margin of 45~deg leaves room for unmodelled lag.
The delay margin covers the case that the estimated delay is too short.
The telemetry latency adds to the delay of the rotor-speed path.
Its value is not part of the evidence, so the delay margin is compared with the control periods only.
It is large in every cell, so this does not change the verdict.

Two reading rules apply.
\begin{enumerate}
  \item If the phase never reaches $-180^\circ$ in the measured band, there is no gain margin, and no gain limit applies.
  \item A crossing of $-180^\circ$ with $|L|>1$ is a gain-reduction margin $20\log_{10}|L|$.
    It must also be at least 6~dB.
\end{enumerate}

\section{Linear models for the cross-check}\label{sec:an_models}

Linear models of the loops predict the margins.
The prediction is compared with the measurement.
The acceptance rule is that $\Delta\mathrm{PM}$ is within 10~deg and $\Delta\mathrm{GM}$ within 3~dB, or that the difference is explained.

The rate model is a discrete-time loop built from the law of \cref{ch:inner_loop}.
It contains the setpoint lead, the actuator time constant $\tau_a$, the delay and the filters.
The attitude loop sits outside the rate loop.
The attitude loop gain is
\begin{equation}\label{eq:an_att}
  L_{\mathrm{att}} = A_{\mathrm{att}}\,T_{\mathrm{cl}} ,\qquad
  A_{\mathrm{att}} = p\,z^{-1}\,\frac{\ts\,z}{z-1} ,\qquad
  T_{\mathrm{cl}} = \frac{N(s)}{1+L_{\mathrm{rate}}} ,
\end{equation}
with the attitude gain $p$ (8 in the reference airframe) and a delay of one sample.
$T_{\mathrm{cl}}$ is the closed rate loop from rate setpoint to rate.
The numerator $N$ is scaled so that $T_{\mathrm{cl}}$ has unit gain at low frequency.
The rate and attitude loops are related by
\begin{equation}\label{eq:an_ident}
  1 + L_{\mathrm{M3a}} = \bigl(1+L_{\mathrm{rate}}\bigr)\bigl(1+L_{\mathrm{att}}\bigr) ,
\end{equation}
where $L_{\mathrm{M3a}}$ is the loop broken at the torque setpoint.
This identity holds to $10^{-9}$ in the model test.

The attitude loop was flown with the velocity and position loops closed.
The model must include them.
The outer-loop factor is
\begin{equation}\label{eq:an_outer}
  O(s) = 1 + \Bigl(K_v + \frac{K_i}{s}\Bigr)\frac{K_p + s}{s^{2}} + K_d\,F(s) ,
\end{equation}
where $F$ is a first-order low-pass at 5~Hz, and $K_p$, $K_v$, $K_i$, $K_d$ are the position and velocity gains.
The attitude model with the outer loops is $L_{\mathrm{att}}\,O$.
The acceleration increment of the INDI stage adds no small-signal path.
No linear model of the break at the input or output of the INDI stage was built.

\section{Results}\label{sec:an_results}

\subsection{Outer loops}\label{sec:an_outer}

The outer-loop cells cover three axes (N, E, D), both break points, two controller rates (400 and 800~Hz) and two seeds.
This gives 24 cells.
\Cref{tab:an_outer} gives the range over the four cells of each group.

\begin{table}[tbp]
\caption{Outer-loop margins of the final configuration, range over four cells (two rates, two seeds) per row. Break point \texttt{pt0} is the input and \texttt{pt1} the output of the INDI stage. ``none'' means no $-180^\circ$ crossing in 0.2 to 8~Hz.}\label{tab:an_outer}
\centering
\small
\begin{tabularx}{\textwidth}{@{}L{2.2cm} L{1.9cm} L{2.8cm} L{2.8cm} Y@{}}
\toprule
Break, axis & GM (dB) & PM (deg) & DM (ms) & Note \\
\midrule
\texttt{pt0}, N & 13.5 to 14.0 & 49.7 to 55.3 & 367 to 399 & \\
\texttt{pt0}, E & 13.5 to 14.0 & 49.1 to 53.0 & 367 to 405 & lowest PM \\
\texttt{pt0}, D & none & 65.4 to 68.2 & 246 to 283 & \\
\texttt{pt1}, N & 13.8 to 14.4 & 52.5 to 71.3 & 464 to 591 & one cell blocked \\
\texttt{pt1}, E & 13.9 to 14.6 & 52.9 to 56.9 & 461 to 485 & one cell re-flown \\
\texttt{pt1}, D & none; 22.4 in one cell & 58.7 to 66.2 & 229 to 288 & two cells blocked \\
\bottomrule
\end{tabularx}
\end{table}

All 24 cells meet the limits.
The smallest phase margin is 49.1~deg, 4.1~deg above the limit.
The smallest gain margin in the horizontal axes is 13.5~dB.
The smallest plateau gain margin is 12.9~dB; it is the second key's reading of the flat gain region.
The smallest delay margin is 229~ms.
At 800~Hz one control period is 1.25~ms, so the delay margin is far above the limit.
The vertical axis has no gain margin in 7 of 8 cells, because the phase stays above $-180^\circ$.
One cell shows 22.4~dB (17.8~dB in key~2).

One cell had a hole in the log.
The controller topic had a single gap of 22.5~ms against a limit of 10~ms excess.
The cell was flown again.
The new run gave PM 56.9~deg, GM 13.9~dB and a clean log.
The table uses the new run.

Three cells are blocked on key agreement.
Both keys are above the limits in each of them.
\begin{itemize}
  \item North, \texttt{pt1}, 800~Hz, seed 1: PM 71.3 and 67.8~deg. The difference exceeds 3~deg.
  \item Down, \texttt{pt1}, 400~Hz, seed 2: PM 58.7 and 58.2~deg. Key~1 has no gain margin, key~2 finds 16.3~dB.
  \item Down, \texttt{pt1}, 800~Hz, seed 2: GM 22.4 and 17.8~dB. The difference exceeds 1~dB.
\end{itemize}
They were not flown again.

\subsection{Rate and attitude loops}\label{sec:an_rate_att}

On the final configuration, one cell was scored for each of two loops.
\begin{itemize}
  \item Rate loop, roll, 800~Hz: PM 46.9~deg, GM 15.7~dB, DM 75.5~ms.
  \item Attitude loop, pitch, 400~Hz: PM 48.0~deg, GM 13.7~dB at 4.29~Hz, DM 93.5~ms.
\end{itemize}
\Cref{tab:an_spot} lists them.
The phase margins are 1.9 and 3.0~deg above the limit.
This is a thin reserve and the evidence is one run each.
Pitch and yaw of the rate loop, and yaw of the attitude loop, were not measured on the final configuration.

\begin{table}[tbp]
\caption{Spot cells of the final configuration. Both have a clean log. The limits are GM 6~dB, PM 45~deg.}\label{tab:an_spot}
\centering
\small
\begin{tabularx}{\textwidth}{@{}L{4.2cm} L{1.7cm} L{1.7cm} L{1.7cm} Y@{}}
\toprule
Cell & GM (dB) & PM (deg) & DM (ms) & Verdict \\
\midrule
Rate, roll, 800~Hz & 15.7 & 46.9 & 75.5 & pass \\
Attitude, pitch, 400~Hz & 13.7 & 48.0 & 93.5 & pass \\
\bottomrule
\end{tabularx}
\end{table}

The attitude cell has two $-180^\circ$ crossings.
One is at 4.29~Hz with $|L|<1$ and a margin of 13.7~dB.
The other is at 0.237~Hz with $|L|=7.3$.
The scorer reports the signed minimum, $-17.3$~dB, for the second one.
Under the reading rule of \cref{sec:an_criteria}, it is a gain-reduction margin of $+17.3$~dB.
The cell passes with GM 13.7~dB.
Read literally as a signed minimum, the cell would fail.
This reading is the one recorded in the verdict.

\subsection{Margin figure}\label{sec:an_fig}

\Cref{fig:g2_margins} shows every flown cell.
Filled markers are 400~Hz and open markers 800~Hz.
Dashed lines are the limits.
Cells without a gain margin are left out of the gain panel.
The delay margin is far above two control periods everywhere.
\indifigure{g2_margins}{Stability margins measured in flight at the final configuration. Each marker is one flown cell; dashed lines are the acceptance thresholds. Loops without a $-180^\circ$ crossing in the measured band have no gain margin. The delay margin is far above two control periods in every cell. The 250~Hz controller rate was not measured.}{fig:g2_margins}

\Cref{fig:bode_rate} (\cref{ch:inner_loop}) shows the measured response of the roll rate cell.
In the Bode plot the gain falls through 0~dB at 1.72~Hz.
The phase there is $-133^\circ$, so the phase margin is 46.9~deg.
The phase reaches $-180^\circ$ at 18.3~Hz, where the gain is $-15.7$~dB.
The Nichols chart shows the same two points on one curve.
The curve passes the critical point $(-180^\circ,\,0~\mathrm{dB})$ with a clear gap.
The plot has 18 measured bins and no model overlay.

\section{Agreement of model and measurement}\label{sec:an_agree}

The cross-checks below use pre-final binaries.
They test the models, not the final margins.

\paragraph{Rate loop.}
The model reproduces an earlier measurement (phase margin 23.6 against 22.3~deg, gain margin 4.7 against 4.1~dB).
It also reproduces the shift in phase margin when the setpoint lead was added (a drop of 9.7~deg in the model, 8.4~deg measured).
Two differences remain.
The model gain is 0.2 to 1.7~dB too high.
The model phase leads the measurement by 17~deg at 1~Hz, and the lead falls to $-5^\circ$ at 8~Hz.
At the crossover the phase margins agree within 1.1~deg.
No gain-margin comparison on the measured path exists.

\paragraph{Attitude loop.}
\Cref{tab:an_att} compares model and flight for the roll cells.
Without the outer loops, the model phase margin was about 15~deg too high.
With the outer-loop factor of \cref{eq:an_outer}, the phase margins agree within 0.3~deg at all three rates.
The gain margin of the model is 1.6 to 4.6~dB lower than measured.
The model is therefore conservative.
The cells at 400 and 267~Hz are outside the 3~dB rule.
The likely cause is the phase error of the rate model near $-180^\circ$, but this was not shown.
Both gain margins are far above 6~dB.

\begin{table}[tbp]
\caption{Attitude loop, roll: model with outer loops against flight. Pre-final binaries.}\label{tab:an_att}
\centering
\small
\begin{tabularx}{\textwidth}{@{}L{2.6cm} L{2.6cm} L{2.6cm} L{2.2cm} Y@{}}
\toprule
Controller rate & Model PM / GM & Measured PM / GM & $\Delta$PM / $\Delta$GM & 10~deg / 3~dB rule \\
\midrule
800~Hz & 57.6 / 20.3 & 57.3 / 21.9 & $-0.3$ / $+1.6$ & pass \\
400~Hz & 56.4 / 18.6 & 56.6 / 22.0 & $+0.2$ / $+3.4$ & GM outside \\
267~Hz & 55.3 / 17.2 & 55.5 / 21.8 & $+0.2$ / $+4.6$ & GM outside \\
\bottomrule
\end{tabularx}
\end{table}

\paragraph{Limits of the cross-check.}
The yaw loops are not modelled.
Pitch is identical to roll in the model, so the model cannot explain a difference between them.
The model runs the attitude controller at the rate-loop rate with a delay of one sample.
The module runs at the attitude callback rate.
The effective velocity and position gains of the flown runs were not read back.
For the outer loops, no model of the break at the stage input or output exists.
A check at three frequencies with the stock horizontal gains gave $|L|$ of 0.60, 0.29 and 0.14 at 0.5, 1 and 2~Hz.
The measurement gave 0.69, 0.27 and 0.16.
This check was not scored by a second key.

\section{Verdict}\label{sec:an_verdict}

The stability gate (G2) is \passed{} on the final configuration.
The recorded result is as follows.
\begin{itemize}
  \item All 24 outer-loop cells meet the limits: PM at least 49.1~deg, horizontal GM at least 13.5~dB, plateau GM at least 12.9~dB, DM at least 229~ms. All logs are clean after one re-flight.
  \item The two spot cells meet the limits (rate PM 46.9~deg, attitude PM 48.0~deg).
  \item The attitude step test passes in 9 of 9 flights with both keys.
\end{itemize}
Four limits of the evidence are recorded.
\begin{enumerate}
  \item Three cells are blocked on key agreement, with both keys above the limits.
  \item The phase margins of the two spot cells are thin, 1.9 and 3.0~deg above the limit, from one run each.
  \item The pitch-attitude cell has a signed gain minimum of $-17.3$~dB at 0.237~Hz. It passes as a gain-reduction margin.
  \item The 250~Hz rate was not measured.
\end{enumerate}
\Cref{ch:validation} reports the flight tests of the other gates.
